%\documentclass{Math_Note}
%\setlength{\parindent}{2em}

%\title{\textbf{Mathmatical Analysis Zorich Comprehensions Chapter 5}}
%\author{Buce-Ithon}
%\newdateformat{mydate}{\twodigit{\THEDAY}{ }\shortmonthname[\THEMONTH], \THEYEAR}
%\date{\today}

%\begin{document}
% Title page
%\maketitle

% Content
%\newpage
%\tableofcontents
%\newpage

% \setcounter{section}{-1}
\newpage
\section{Chapter 5.Differential Calculus}
\subsection{Comprehensions}
1st part: Introduce the definition of functions differentiable at a point and related concepts: differential, derivative and increment. 
Besides, it also includes the geometrical meaning of them and some calculation examples. 
After basic definitions, we are given the basic rules of differentiation, including: arithmetic operations, differentiation of composite 
functions, differentiation of inverse functions, and higer order derivatives (Lebniz's formula). And also derivatives of basic elementary 
functions. \\
2nd part: Then we dive into some basic theorems of differential calculus: /All of beging-Most basic/ \underline{Fermat's Lemma}, /Dev on 
basic-3 Mean Value Theorems/ \underline{Rolle's Theorem} $\to$ \underline{Lagrange's Mean Value Theorem} $\to$ \underline{Cauchy's Mean 
Value} \underline{Theorem} (due to last 2 of them associated the increment and derivative, they're also called Theorem of Lagrange and 
Cauchy on finite increments). Finally, we're given the most important theorem in differential calculus: \underline{Taylor's Formula} (the 
reminder term could be obtained using Cauchy's Mean Value Theorem or Theorem of Cauchy on finite increments). \\
3nd part: It's the last part of this chapter, which focuses on the applications (analyze properties of functions) of differential calculus, 
including: monotoncity, interior extremum, convacity, and graph sketching. And the well-known L'H\^opital's Rule. Remains lots of examples 
in Problems on Natural Science. \\
DLC: The special (which could be called overall) point is that Zorich left a specific section for complex numbers and functions, even including 
analyticity (or holomorphism). 
\newline\newline
\fbox{\begin{minipage}{0.95\textwidth}
\textbf{Definition: } /differential/ $\underbrace{\underbrace{f(x+h)-f(x)}=\overbrace{A(x)\underbrace{h}}^{\text{differential}}}_{\text{increments}}+
\underbrace{o(a;h)}_{o(h),\ h\to 0}$, $x \in E$ 
/derivative/ $f'(x) = A(x) = \lim_{h\to 0}\frac{f(x+h)-f(x)}{h}$
\end{minipage}}
\newline\newline
All of beging-Most basic \\
\fbox{\begin{minipage}{0.95\textwidth}
\begin{thm}
/Fermat's Lemma/ \\
$f : E \to \mathbb{R}$ is differentiable at an interior extremum $x_0 \in E$ $\implies$  $f'(x_0) = 0$.
\end{thm}
\end{minipage}}
\newline\newline
3 Mean Value Theorems \\
\fbox{\begin{minipage}{0.95\textwidth}
\begin{thm}
/Rolle's theorem/ \\
$f : [a, b] \to \mathbb{R}$ is continuous on $[a, b]$ and differentiable on $(a, b)$ and $f(a) = f(b)$, then $\exists \xi \in (a, b)$ s.t. $f'(\xi) = 0$.
\end{thm}
\end{minipage}}
\newline\newline
\fbox{\begin{minipage}{0.95\textwidth}
\begin{thm}
/Lagrange's Mean Value Theorem/ \\
$f : [a, b] \to \mathbb{R}$ is continuous on $[a, b]$ and differentiable on $(a, b)$, then $\exists \xi \in (a, b)$ s.t. $f(b) - f(a) = f'(\xi)(b - a)$.
\end{thm}
\end{minipage}}
\newline
\fbox{\begin{minipage}{0.95\textwidth}
\begin{thm}
/Cauchy's Mean Value Theorem/ \\
Let $x = x(t)$ and $y = y(t)$ be continuous on $[\alpha, \beta]$ and differentiable on $(\alpha, \beta)$.
Then $\exists \tau \in [\alpha, \beta]$ s.t. $x'(\tau)(y(\beta) - y(\alpha)) = y'(\tau)(x(\beta) - x(\alpha))$. \\
If in addition $x'(t) \neq 0$, $\forall t \in [\alpha, \beta]$, then $x(\alpha) \neq x(\beta)$ and $\frac{y(\beta) - y(\alpha)}{x(\beta) - x(\alpha)} 
= \frac{y'(\tau)}{x'(\tau)}$.
\end{thm}
\end{minipage}}
\newline\newline
Master in differential calculus \\
\fbox{\begin{minipage}{0.95\textwidth}
\begin{thm}
/Taylor's Formula/ \\
$$F(a) = f(x_0) + f'(x_0)(x - x_0) + \cdots + \frac{f^{(n)}(x_0)}{n!}(x - x_0)^n + r_n(x_0, x)$$
\end{thm}
\end{minipage}}
\newline\newline
\fbox{\begin{minipage}{0.95\textwidth}
\textbf{Remark: } /Reminder term/ \\
$$r_n(x_0;x) = \frac{\varphi(x) - \varphi(x_0)}{\varphi'(\xi)n!} f^{(n+1)}(\xi)(x-\xi)^n.$$
or $$\frac{r_n(x_0;x)}{\varphi(x) - \varphi(x_0)}=\frac{f^{(n+1)}(\xi)(x-\xi)^n}{\varphi'(\xi)n!} \ \ \text{/Cauchy's MVT/}$$
/Cauchy’s formula of the remainder term/ \\
$$r_n(x_0; x) = \frac{1}{n!} f^{(n+1)}(\xi)(x - \xi)^n(x - x_0).$$
A particularly elegant formula results if we set $\varphi(t) = (x - t)^{n+1}$ \\
/The Lagrange form of the remainder/ \\
$$r_n(x_0; x) = \frac{1}{(n + 1)!} f^{(n+1)}(\xi)(x - x_0)^{n+1}.$$
We remark that when $x_0 = 0$ Taylor's formula is often called MacLaurin's formula.
\end{minipage}}

\newpage
\subsection{Problems Sets}
\begin{prb}
    (C5S1Q5) /Continuous, but nowhere derivative/ \\
    Set $\Psi_0(x) = \begin{cases} x, & \text{if } 0 \leq x \leq \frac{1}{2}, \\ 1 - x, & \text{if } \frac{1}{2} \leq x \leq 1, \end{cases}$ and extend this 
    function to the entire real line so as to have period 1. We denote the extended function by $\varphi_0$. Further, let $\varphi_n(x) = \frac{1}{4^n} 
    \varphi_0(4^n x)$. The function $\varphi_n$ has period $4^{-n}$ and a derivative equal to +1 or $-1$ everywhere except at the points $x = \frac{k}{2^{n+1}}, 
    n \in \mathbb{Z}$. Let $f(x) = \sum_{n=1}^\infty \varphi_n(x)$. \\
    Show that the function $f$ is defined and continuous on $\mathbb{R}$, but does not have a derivative at any point.
    (This example is due to the well-known Dutch mathematician B. L. van der Waerden (1903–1996). The first examples of continuous functions having no derivatives 
    were constructed by Bolzano (1830) and Weierstrass (1860).)
\end{prb}
\fbox{\begin{minipage}{0.95\textwidth}
\begin{sol*}
1. Using Weierstrass M-test (uniform convergence of series of functions) and the uniform limit theorem (the uniform limit of any sequence of continuous functions 
is continuous), we can show that $f \in C(\mathbb{R})$. \\
2. Just consider $\lim_{n \to \infty} \frac{f(x\pm 4^{-n})-f(x)}{4^{-n}}=\lim_{n \to \infty} \frac{\sum_{k=0}^{\infty}\varphi_k(x\pm 4^{-n})-\sum_{k}\varphi_k(x)}
{4^{-n}}$, and using the periodicity of $\varphi_k$.  
\end{sol*}
\end{minipage}}
\newline\newline
\fbox{\begin{minipage}{0.95\textwidth}
\textcolor{cyan}{/Dev/} \textbf{More step problem:} \\
Prove that the nowhere differential functions are dense in $C(\mathbb{R})$.
\end{minipage}}

\begin{prb}
    (C5S2Q1) \\
    Let $\alpha_0, \alpha_1, \dots , \alpha_n$ be given real numbers. Find a polynomial $P_n(x)$ of degree $n$ having the derivatives $P_n^{(k)}(x_0) = \alpha_k, 
    k = 0, 1, \dots , n$, at a given point $x_0 \in \mathbb{R}$.
\end{prb}
\fbox{\begin{minipage}{0.95\textwidth}
\begin{qsol*}
$P_n(x)=\sum_{k=0}^{n} \alpha_k(x-x_0)^k$ 
\end{qsol*}
\end{minipage}}

\begin{prb}
    (C5S2Q2) 
\end{prb}
\fbox{\begin{minipage}{0.95\textwidth}
\begin{qsol*}
\begin{enumerate}
    \item[a)] $f(x)=\begin{cases} e^{-\frac{1}{x^2}} & x\neq 0 \\ 0 & x=0 \end{cases}$,$f'(x)=\begin{cases} \frac{2e^{-\frac{1}{x^2}}}{x^3} 
        & x\neq 0 \\ 0 & x=0 \end{cases}$, $f^{(n)}(x)=\lim_{x\to\infty}\frac{f^{(n-1)}(x)-0}{x-0}=[\text{Induction}]=\lim_{x\to\infty}\frac{R_n(x)}
        {e^{\frac{1}{x^2}}}=0$. 
    \item[b)] $f(x)=\begin{cases} x^2\sin(\frac{1}{x}) & x\neq0 \\ 0 & x=0 \end{cases}$, $f(x)=\begin{cases} 2x\sin(\frac{1}{x})-\cos(\frac{1}{x}) & x\neq0 
        \\ 0 & x=0 \end{cases}$ 
    \item[c)] $f(x)=\begin{cases} e^{-\frac{1}{(1+x)^2}}\cdot e^{-\frac{1}{(1-x)^2}} & -1<x<1 \\ 0 & 1\leq|x| \end{cases}$ 
\end{enumerate}
\end{qsol*}
\end{minipage}}

\begin{prb}
    (C5S2Q3) \\
    $$\left[ {x^{n - 1} f\left( {\frac{1}{x}} \right)} \right]^{(n)} = \frac{{( - 1)^n }}{{x^{n + 1} }}f^{(n)} \left( {\frac{1}{x}} \right)$$
\end{prb}
\fbox{\begin{minipage}{0.95\textwidth}
\begin{qsol*}
$[$Induction$]$. Assume that the base case $n=1$ is proved. \\
Assuming that $\frac{d^n}{dx^n}\left[x^{n-1}f\left(\frac{1}{x}\right)\right]=\frac{(-1)^n}{x^{n+1}}f^{(n)}\left(\frac{1}{x}\right)$, we have $$\frac{d^{n+1}}
{dx^{n+1}}\left[x^{n}f\left(\frac{1}{x}\right)\right]=\frac{d^n}{dx^n}\left[nx^{n-1}f\left(\frac{1}{x}\right)-x^{n-2}f'\left(\frac{1}{x}\right)\right]$$ $$
\underbrace{=\frac{n(-1)^{n}}{n^{n+1}}f^{(n)}\left(\frac{1}{x}\right)}_\text{by asumption}-\frac{d}{dx}\left(\frac{d^{n-1}}{dx^{n-1}}\left[x^{n-2}f'\left(
\frac{1}{x}\right)\right]\right)$$ $$=\frac{n(-1)^{n}}{n^{n+1}}f^{(n)}\left(\frac{1}{x}\right)-\frac{d}{dx}\underbrace{\left[\frac{(-1)^{n-1}}{x^n}f^{(n)}\left(
\frac{1}{x}\right)\right]}_\text{by assumption}$$ $$=\frac{n(-1)^{n}}{n^{n+1}}f^{(n)}\left(\frac{1}{x}\right)+\frac{n(-1)^{n-1}}{x^{n+1}}f^{(n)}\left(\frac{1}{x}\right)
+\frac{(-1)^{n+1}}{x^{n+2}}f^{(n+1)}\left(\frac{1}{x}\right)$$ $$=\frac{(-1)^{n+1}}{x^{n+2}}f^{(n+1)}\left(\frac{1}{x}\right)\square$$
And that about wraps it up. \textit{"from AOPS"}
\end{qsol*}
\end{minipage}}

\begin{prb}
    (C5S2Q4) \\
    Let $f$ be a differentiable function on $\mathbb{R}$. Show that
    \begin{enumerate}
        \item[a)] if $f$ is an even function, then $f'$ is an odd function;
        \item[b)] if $f$ is an odd function, then $f'$ is an even function;
        \item[b)] ($f'$ is odd) $\Leftrightarrow$ ($f$ is even).
    \end{enumerate}
\end{prb}
\fbox{\begin{minipage}{0.95\textwidth}
\begin{qsol*}
By definition.
\end{qsol*}
\end{minipage}}

\begin{prb}
    (C5S2Q5) \\
    Show that
    \begin{enumerate}
        \item[a)] the function $f(x)$ is differentiable at the point $x_0$ if and only if $f(x) - f(x_0) = \varphi(x)(x - x_0)$, where $\varphi(x)$ is a function that 
        is continuous at $x_0$ (and in that case $\varphi(x_0) = f'(x_0)$);
        \item[b)] if $f(x) - f(x_0) = \varphi(x)(x - x_0)$ and $\varphi \in C^{(n-1)} \left( U(x_0) \right)$, where $U(x_0)$ is a neighborhood of $x_0$, then $f(x)$ has a 
        derivative $(f^{(n)}(x_0))$ of order $n$
    \end{enumerate}
\end{prb}
\fbox{\begin{minipage}{0.95\textwidth}
\begin{qsol*}
By definition.
\end{qsol*}
\end{minipage}}

\begin{prb}
    (C5S2Q6) \\
    Give an example showing that the assumption that $f^{-1}$ be continuous at $y_0$ cannot be omitted from following theorem: \\
    \begin{thm}
        (The derivative of an inverse function). Let the functions $f : X \to Y$ and $f^{-1} : Y \to X$ be mutually inverse and continuous at points $x_0 \in X$ and 
        $f(x_0) = y_0 \in Y$ respectively. If $f$ is differentiable at $x_0$ and $f'(x_0) \neq 0$, then $f^{-1}$ is also differentiable at the point $y_0$, and 
        $(f^{-1})'(y_0) = (f'(x_0))^{-1}$.
    \end{thm}
\end{prb}
\fbox{\begin{minipage}{0.95\textwidth}
\begin{qsol*}
Let $f(x)=\begin{cases} x+x^2\sin\frac{1}{x}, & x\neq 0 \\ 0, & x=0 \end{cases}$, after calculating, we could find $f$ is continuous and differentiable, but $f^{-1}$ 
doesn't even exist. 
\end{qsol*}
\end{minipage}}

\begin{prb}
    (C5S3Q1) \\
\end{prb}

%\end{document}